\chapter{State of the Art}
\label{ch:ch2_soa}

% [DRAFT v1] Source material: docs/estado_del_arte_trng.md, estado_del_arte_cripto.md,
% plataforma_y_flujo.md, via_asic.md. Every numeric claim below carries a
% reference that was checked against its source when the review was compiled;
% items the review could not open are marked as such in the working notes.

The review is organized along the chain the thesis builds: how ring
oscillators are turned into random bits, how their entropy is modeled and
bounded, how their jitter is measured from inside the device, what the two
governing standards demand, which symmetric primitives fit a constrained
engine, what the commercial parts offer, and what changes when the design
leaves the FPGA for silicon. The last section states the gap this work
addresses, in the guarded form the working rules require.

\section{Ring-oscillator TRNG architectures}
\label{sec:soa_archs}

A ring oscillator is an odd chain of inverters closed on itself; its
period fluctuates because each transition is delayed by a random amount,
and that fluctuation --- jitter --- is the noise source of an entire family
of generators. What distinguishes the architectures is how the jitter is
accumulated and sampled~\cite{petura_2016_survey, lubicz_2024_recommendations}.

The \textbf{elementary ring oscillator} (ERO) samples one ring with a
flip-flop clocked by a second ring divided by $K_D$; the divider sets the
accumulation time. It has the most mature stochastic
model~\cite{baudet_2011_security, killmann_2008_design, ma_2014_entropy} and
an embedded jitter measurement~\cite{fischer_2014_embedded}, at the price
of a very low output rate ($K_D$ of $10^4$ to $10^5$: 80\,000 in Spartan-6
and 135\,000 in Cyclone~V in Petura's implementation, about 430\,000 in
Fischer and Lubicz's design for 0.997~bit of entropy). The
\textbf{multi-ring} generator of Sunar, Martin and
Stinson~\cite{sunar_2007_provably} XORs many rings sampled by a reference
clock; Wold and Tan~\cite{wold_2009_analysis} added a flip-flop per ring
before the XOR and reported 100~Mbit/s in under 100 logic elements. Its
weakness was exposed by Bochard \emph{et al.}~\cite{bochard_2010_true}:
part of the apparent randomness is \emph{pseudo}-randomness from the
interaction of rings at different frequencies, and the independence the
proof assumes does not hold on an FPGA. \textbf{Coherent sampling} (COSO)
uses two rings of nearly equal period and counts the beat~\cite{kohlbrenner_2004_embedded,
valtchanov_2009_enhanced}; its model gives min-entropy directly from the
distribution of the count~\cite{peetermans_2021_configurable}, and
dynamic calibration of reconfigurable rings makes it portable across
devices and placements (3.30~Mbit/s in Spartan-6, 4.65~Mbit/s in
Spartan-7)~\cite{peetermans_2019_portable, peetermans_2021_configurable}.
The \textbf{transition-effect ring} (TERO) counts how many times a
bistable structure oscillates before settling~\cite{varchola_2010_new},
with a full physical-to-stochastic model~\cite{bernard_2019_physical};
the \textbf{self-timed ring} (STR) circulates several events in an
asynchronous ring to obtain phases spaced by the order of the
jitter~\cite{cherkaoui_2013_very}, reaching 154~Mbit/s; and the
\textbf{ES-TRNG} of Yang \emph{et al.} uses the carry chain as a
time-to-digital converter to encode the edge position with 10~LUTs and
5~flip-flops~\cite{yang_2018_estrng, rozic_2015_highly}.

Table~\ref{tab:petura} reproduces the comparison Petura \emph{et al.}
made of the AIS-compliant cores on one device~\cite{petura_2016_survey}.
The entropy column is a statistical estimate (the AIS~31 T8 test), not a
model bound; the feasibility column is the authors' qualitative note,
where COSO and TERO score worst because they need manual adjustment per
device. No comparable table exists for Artix-7; the open OpenTRNG project
implements ERO, multi-ring and COSO on an Arty~A7 with the same XC7A35T
die as the Basys~3 but publishes neither entropy nor throughput
figures~\cite{opentrng_2025_docs}.

\begin{table}[htbp]
\centering
\caption{AIS-compatible TRNG cores on Xilinx Spartan-6 (45~nm), after
Petura \emph{et al.}~\cite{petura_2016_survey}, Table~II. Entropy is the
T8 estimate; ``feasibility'' is the authors' rank (1 best).}
\label{tab:petura}
\begin{tabular}{@{}lrrrrl@{}}
\toprule
Core & LUT/reg. & Power (mW) & Rate (Mbit/s) & Entropy/bit & Feasib. \\
\midrule
ERO  & 46/19   & 2.16  & 0.0042 & 0.999 & 5 \\
COSO & 18/3    & 1.22  & 0.54   & 0.999 & 1 \\
MURO & 521/131 & 54.72 & 2.57   & 0.999 & 4 \\
PLL  & 34/14   & 10.6  & 0.44   & 0.981 & 3 \\
TERO & 39/12   & 3.31  & 0.625  & 0.999 & 1 \\
STR  & 346/256 & 65.9  & 154    & 0.998 & 2 \\
\bottomrule
\end{tabular}
\end{table}

\section{Stochastic models and entropy bounds}
\label{sec:soa_models}

The requirement that a physical generator come with a \emph{stochastic
model} --- a mathematical description of the noise source from which the
entropy of the output can be derived --- was introduced by Killmann and
Schindler~\cite{killmann_2008_design, schindler_2002_evaluation} and is
mandatory in AIS~20/31. For ring oscillators the reference model is
Baudet \emph{et al.}~\cite{baudet_2011_security}: the phase of the sampled
ring is a Wiener process with drift $\mu$ and volatility $\sigma^2$, the
sampled bit is the half-period the phase falls in, and everything depends
on the quality factor $Q=\sigma^2\Delta t$ (accumulated phase variance
between samples) and the frequency ratio $\nu=\mu\Delta t$. The model
gives the bit probability conditioned on the previous phase, the
probability of any output vector, and a lower bound on the Shannon
entropy rate, $H \ge 1 - \frac{4}{\pi^2\ln 2}e^{-4\pi^2 Q}$, whose exact
status is examined in Chapter~\ref{ch:ch3_entropy}. Petura \emph{et al.}
rewrite it in design quantities, $H = 1 - \frac{4}{\pi^2\ln2}\exp(-\pi^2
\sigma_{th}^2 K T_2/T_1^3)$, and Fischer and Lubicz invert it to obtain
the divider $K_D$ from a measured jitter~\cite{fischer_2014_embedded}. Ma
\emph{et al.} give an alternative model used by the Leuven
group~\cite{ma_2014_entropy}; Haddad \emph{et al.} question the mutual
independence of jitter realizations on which these cancellations
rest~\cite{haddad_2014_assumption}; and Saarinen warns that Baudet's bound
``is never lower than 0.415 even when $Q$ approaches zero'' and is safe
only under additional assumptions, proposing a bound valid over the whole
range~\cite{saarinen_2021_entropy}.

Two refinements matter for a design. First, not all jitter is entropy.
Fischer and Lubicz separate the accumulated phase variance into a
random-walk term, linear in time, produced by thermal (white) noise, and
flicker ($1/f^\beta$) terms whose variance grows quadratically and
dominate at long accumulation times~\cite{fischer_2014_embedded};
deterministic, global jitter from the supply and substrate behaves like
the latter and can be known or induced by an
adversary~\cite{fischer_2008_enhancing}. Only the linear term may be
credited, which is why a divider cannot be made arbitrarily large to
compensate a weak source. Second, every embedded measurement is
\emph{differential}: two rings sharing a supply see the same global
jitter, which cancels in their relative phase, leaving the sum of their
local jitters. Lubicz and Sk\'orski formalize this ``jitter transfer
principle'' and show that with two rings one obtains only the sum, while
three rings and three pairwise measurements recover the individual
variances; their experiment also shows the assumption of variance linear
in period failing by a factor of two in one case~\cite{lubicz_2024_jitter}.

\section{Embedded jitter measurement}
\label{sec:soa_jitter}

Oscilloscopes cannot resolve the picoseconds of jitter of a ring inside a
chip without an external divider~\cite{benea_2024_flicker}, so the useful
methods measure from within. The \textbf{counter method} counts periods of
one ring in a fixed window of the other; the variance of the count grows
with the accumulated relative jitter~\cite{killmann_2008_design,
lubicz_2015_towards}, with a quantization floor of one count and
contamination by flicker at long windows. The method of \textbf{Fischer
and Lubicz}~\cite{fischer_2014_embedded} needs no divider: in $K\approx10^4$
blocks of $N\approx100$ bits sampled at $K_D=1$, the fraction of pairs
$(b_j,b_{j+M})$ that differ tracks the phase accumulated over $M$ periods,
and its variance across blocks, plotted against $M$ between 200 and 1600,
gives the jitter per period with an error under 5\,\% in simulation; on a
Cyclone~III they measured $\sigma=5.01$~ps per 7.81~ns period, and showed
that cooling the chip reduces $\sigma$ --- which makes the same measurement
an online test. \textbf{Coherent sampling} yields the accumulated jitter
from the variance of the beat count with no extra
hardware~\cite{peetermans_2021_configurable}. Most recently Benea
\emph{et al.} characterize the jitter of rings with the \textbf{Allan
variance}, fitting $\sigma_t^2(t)=a_0+a_1 t+a_2 t^2$ to separate
quantization, thermal and flicker terms, and do so both on an Artix-7
--- the FPGA family of this work --- and on a 28~nm FD-SOI circuit, with
data and scripts released~\cite{benea_2024_flicker}. Published per-period
jitter figures cluster at 2--4~ps for periods of 4--8~ns in Spartan-6 and
Cyclone~V~\cite{petura_2016_survey}, 5~ps in Cyclone~III measured
embedded~\cite{fischer_2014_embedded}, and 2--2.5~ps in Cyclone~III and
Virtex-5 for self-timed rings~\cite{cherkaoui_2013_very}; the Artix-7
coefficients of Benea \emph{et al.} are published in counter units, not
seconds, so an absolute figure for the XC7A35T remains to be measured.

\section{Standards: AIS 20/31 v3.0 and NIST SP 800-90B}
\label{sec:soa_standards}

Two documents govern how a physical generator is assessed. The German
\textbf{AIS~20/31}, in its version~3.0 of September 2024~\cite{peter_2024_ais31},
replaces the 2011 requirement of 0.997~bit of Shannon entropy per
internal bit~\cite{killmann_2011_ais31} with a choice in class PTG.2:
Shannon entropy $\ge0.9998$ and/or min-entropy $\ge0.98$, together with
$P(1)\in(0.493, 0.507)$; the BSI's transition policy allows new
evaluations under the old version only until March
2026~\cite{bsi_2025_transition}. Its requirements are structural as much
as numerical: a verifiable stochastic model is mandatory; the online test
must be derived from the model, failing with high probability when the
source parameter leaves the admissible set; a total-failure test must
detect entropy collapsing to zero, and the repetition-count test of
SP~800-90B is explicitly accepted for it; and the black-box suite
$T_{\mathrm{irn}}$ now includes two of NIST's predictors (MultiMMC and
LZ78Y), a deliberate harmonization~\cite{peter_2024_ais31, schindler_2024_ecw}.

The American \textbf{SP~800-90B}~\cite{turan_2018_sp80090b} estimates
\emph{min-entropy} --- the adversarial quantity, $-\log_2$ of the most
probable value --- from at least one million raw samples and a
1000$\times$1000 restart matrix, using the most-common-value estimator
if the source is shown to be i.i.d.\ and otherwise the minimum of ten
estimators, four of them predictors that try to guess the next sample.
It approves two continuous health tests and requires no others: the
Repetition Count Test with cutoff $C=1+\lceil-\log_2\alpha/H\rceil$ and the
Adaptive Proportion Test over a window of 1024 samples for binary sources
(512 for non-binary), $C=1+\mathrm{CRITBINOM}(W,2^{-H},1-\alpha)$, with
$\alpha=2^{-20}$ recommended. It lists the conditioning components whose
output may claim full entropy --- HMAC, CMAC, CBC-MAC, approved hashes and
the two derivation functions of SP~800-90A --- and gives the formula for
the entropy of a conditioned output as a function of input length, input
entropy and internal width. The reference implementation of the
estimators is public~\cite{nist_ea_github}.

The two standards agree on what the older statistical suites cannot do.
SP~800-22~\cite{bassham_2010_sp80022} tests uniformity and independence of
a bit stream; a deterministic generator with zero entropy passes it by
construction, and so does a weak physical source behind a cryptographic
post-processor~\cite{saarinen_2021_entropy}. NIST's 2022 decision to
revise the publication states the point in so many words, ``rejecting its
use for assessing cryptographic random number generators''~\cite{nist_2022_decision}.
The empirical record agrees: a certified smart-card generator showed
consistent biases while passing most tests~\cite{hurleysmith_2018_certifiably,
hurleysmith_2018_great}, and the ESP32's generator passes the same
screenings with its radio off, when Espressif's own documentation says it
returns pseudo-random values~\cite{espressif_idf_random, espressif_trm_esp32,
blancoromero_2026_entropy}. The academic literature on that generator
applies only SP~800-22 or Dieharder~\cite{castillo_2026_automated,
hidayatulloh_2024_performance}; no SP~800-90B characterization of it was
found.

\section{Symmetric cores for constrained hardware}
\label{sec:soa_cores}

The engine needs a block cipher and, if it uses one, a hash. For AES-128
(ten rounds~\cite{nist_2023_fips197}) the design space runs from
fully unrolled pipelines for throughput to 8-bit datapaths for area:
Good and Benaissa report 124 slices and two block RAMs at 2.2~Mbit/s on
a Spartan-II at one extreme and 25~Gbit/s at the
other~\cite{good_2005_fastest}; Chodowiec and Gaj's compact design uses 222
slices and three block RAMs~\cite{chodowiec_2003_compact}; and the
smallest ASIC realizations reach 2.4~kGE with an 8-bit datapath at 226
cycles per block~\cite{moradi2011pushing}, following the composite-field
S-box line of Satoh~\cite{satoh_2001_compact}, Canright~\cite{canright_2005_compact}
and Boyar--Peralta~\cite{boyar_2010_minimization}. The composite-field
constructions optimize ASIC gate counts and lose their advantage in an
FPGA LUT, where a direct table is both faster and simpler. Open
reference designs with published Artix-7 figures exist: a one-round-per-cycle
encryptor at 1104~LUTs and 294~MHz~\cite{hadipour_aesvhdl}, and a 32-bit
iterative encryptor/decryptor at 2298 slices and 97~MHz~\cite{strombergson_aes}.
For an engine whose output leaves through an I2C link, any of these is
far faster than the link, and the choice is dictated by verifiability and
area, not speed.

SHA-256~\cite{nist_2015_fips1804} in its iterative form costs about 750
slices at 108~MHz on Artix-7~\cite{strombergson_sha256}, comparable to a
compact AES; its role in a generator would be as the vetted conditioner of
SP~800-90B. The deterministic generators of
SP~800-90A~\cite{barker_2015_sp80090ar1} --- Hash\_DRBG, HMAC\_DRBG and
CTR\_DRBG --- differ mainly in cost: CTR\_DRBG with AES-128 keeps a state
of 256~bits plus a counter and produces 128~bits per block encryption,
Hash\_DRBG needs 440-bit modular additions, and HMAC\_DRBG two hash
compressions per output block. Without a derivation function CTR\_DRBG
requires full-entropy input of exactly the seed length; with one, the
usual configuration in practice, it accepts arbitrary input. Woodage and
Shumow's analysis warns that the standard's flexibility permits choices
that worsen state leakage and recommends updating after every generate
and reseeding far below the $2^{48}$ limit~\cite{woodage_2019_analysis};
SP~800-90C, final since 2025, frames how a physical source and a DRBG
combine~\cite{barker_2025_sp80090c}. For authentication, AES-CMAC
(SP~800-38B~\cite{nist_2016_sp80038b}, RFC~4493~\cite{song_2006_rfc4493})
costs one block encryption per message block and is itself a vetted
conditioner. The official vectors for all of these are
public~\cite{nist_cavp_blockciphers, nist_cavp_modes, nist_cavp_drbg,
nist_acvp_server}.

\section{Commercial secure elements}
\label{sec:soa_se}

The product category this engine belongs to exists. A secure element is
a small integrated circuit that generates and stores keys and performs
symmetric and asymmetric operations for a host microcontroller, almost
always over I2C. Four parts dominate the low-cost segment: Microchip's
ATECC608B, NXP's SE050, Infineon's OPTIGA Trust~M and ST's STSAFE-A110~\cite{microchip_atecc608b_ds,
nxp_se050_ds, infineon_optiga_trustm_ds, st_stsafe_a110_ds}. The
ATECC608B, at 0.90~\$ in single units~\cite{digikey_prices_2026}, offers
AES-128, SHA-256/HMAC, ECDSA/ECDH on P-256, sixteen key slots and an I2C
interface, and its entropy source holds a NIST validation under
SP~800-90B, obtained precisely with ring oscillators~\cite{nist_esv_e46}.
Two of its competitors make no reference to SP~800-90 in their datasheets;
the third states that its generator is ``designed to meet'' the standard.
The lesson for this work is where an academic design can and cannot
compete: not on price, where a certified part costs less than a euro, but
on the transparency of the entropy, which is exactly what a black-box part
cannot offer and what the standards now demand.

\section{From FPGA to ASIC: what changes in the entropy source}
\label{sec:soa_asic}

Digital logic ports from an FPGA to standard cells with a known penalty
--- Kuon and Rose measured 40$\times$ in area, 3.2$\times$ in delay and
12$\times$ in dynamic power for logic-only designs in 90~nm, and 27--49$\times$
in area for AES specifically~\cite{kuon2006}; the ring oscillator does not
port at all, because in silicon it is built from inverter cells whose
noise, variability and coupling differ from a LUT chain. Three published
findings define what a silicon implementation must add. Benea
\emph{et al.} show that flicker noise, whose weight grows in newer nodes,
changes the character of the jitter and, against expectation, that
treating it as a legitimate source can raise throughput fourfold at equal
entropy~\cite{benea_2024_flicker}. Yang, Blaauw and Sylvester show that
process variation --- larger in 40~nm than in 180~nm --- must be absorbed
by a configurable ring with a tuning loop for the generator to survive
$-40$ to 120~$^\circ$C and 0.6 to 0.9~V~\cite{yang_2016_jssc}. And two
attacks show why a security device cannot ignore its environment: a
frequency injected on the supply locks the rings and removed the jitter
of a commercial secure microcontroller, reducing its key space from
$2^{64}$ to 3300~\cite{markettos_2009_frequency}, and a few hundred
microwatts of radiated field biased a fifty-ring generator until it
failed every test~\cite{bayon_2012_contactless}; the published defence
combines supply filtering with the same calibration loop~\cite{yang_2016_jssc},
and device aging under bias-temperature instability and hot-carrier
injection degrades poorly-timed generators over years while leaving
well-timed ones intact.

Access to silicon is more open than it appears. The University of Sevilla
and IMSE-CNM are members of Europractice~\cite{europractice_memberlist},
whose academic price list is public: a mini@sic block costs 3\,430~\texteuro{}
in UMC~180~nm and 3\,691~\texteuro{} in TSMC~65~nm, and IHP's SG13G2 is
offered at 5\,110~\texteuro{}/mm$^2$ with a free run for open-source designs
under 2~mm$^2$~\cite{europractice_prices2026, europractice_minisic,
ihp_openpdk}. The institute's hardware-security area has already taken
ring-oscillator primitives from Xilinx 7-series FPGAs, evaluated with
SP~800-22 and SP~800-90B, to characterized 65~nm silicon~\cite{imse_area_hwsec,
martinez2022efficientropuf, ortega2024vlsi65nm, ortega2026ropuf65nm}, and
has fabricated cryptographic circuits before~\cite{mora2020trivium}. The
open standard-cell libraries that make a synthesis chapter possible ---
IHP SG13G2, SkyWater SKY130, GlobalFoundries GF180MCU and the predictive
Nangate 45~nm --- come with liberty, LEF and GDS views, and the literature
compares such designs in gate equivalents normalized to the minimum NAND2
of the library~\cite{ihp_openpdk_repo, skywater_pdk_cells, gf180mcu_pdk,
nangate45, kaps2019lwc}, with the caution that logic-synthesis area is a
lower bound that grows by a factor of 1.7 to 4.6 to routed core area.

\section{Gap addressed by this work}
\label{sec:gap}

No absolute novelty is claimed. The ERO and its model, the embedded
jitter measurement, the survey of AIS-compliant cores, the calibrated
COSO on 7-series devices and an open implementation on the very die of
the Basys~3 all exist, and the institute's own group has crossed from
7-series FPGAs to 65~nm silicon with ring-oscillator primitives. What was
\emph{not found} in the literature consulted, and what this work sets out
to provide in a modest but verifiable way, is the following.

\begin{enumerate}
\item A per-period jitter figure for the Artix-7 XC7A35T in absolute
units, obtained inside the device with a differential method, with the
thermal and flicker components separated and with stated uncertainty ---
the published Artix-7 coefficients being in counter units.
\item The complete chain \emph{measured jitter $\to$ quality factor $\to$
min-entropy bound $\to$ divider $\to$ health-test cutoffs}, applied to one
platform and confronted with the ten non-i.i.d.\ estimators of SP~800-90B
and with the PTG.2 requirements of AIS~20/31 \emph{version~3.0}, which
most of the earlier literature could not use.
\item An SP~800-90B characterization of the ESP32 generator in its
operating states, and a controlled comparison between a white-box FPGA
source and a black-box commercial one under the same methodology.
\item The use of a commodity microcontroller with its own cryptographic
accelerators as an independent oracle for an FPGA engine, over a real
bus, with known-answer, interoperability and challenge-response tests.
\item A single-primitive key-generation engine --- AES for conditioning,
generation and authentication --- written as portable RTL, verified
against official vectors, synthesized on an open 130~nm library with the
result stated as a lower bound, and costed for an academic shuttle.
\end{enumerate}

Each item is formulated as ``not found'' rather than ``does not exist''.
One earlier claim of this kind in the working notes --- that no Artix-7
jitter measurement had been published --- had to be corrected when
Benea \emph{et al.}'s paper was found, and the correction is kept visible
there.
